\section{Limitations and conclusion}
\label{sec:conclusion}

\paragraph{Limitations.} The sensing budget uses textbook closed forms:
spherical-then-cylindrical spreading rather than a ray or parabolic-equation
propagation model, a single sea state and shipping level, a generic Wenz
spectrum rather than a measured deep-sea soundscape, and an eDNA plume with
illustrative shedding and decay rates. The synthetic species are four
parametric families; real deep-sea sound includes many mechanisms we did not
model and many animals that make no sound at all, for which acoustics offers
nothing and only cameras, eDNA or trawls apply. The real-data result is a
pilot: \RealClips\ clips of \RealSpecies\ species, a CNN trained from scratch,
three holdout draws; it demonstrates the gap to synthetic performance, not
the ceiling of the method, and the scaled configuration replaces the
from-scratch CNN by a pretrained audio embedding. Clustering uses one
criterion (silhouette) whose failure at larger $U$ we report rather than
tune away. The survey model assumes detections are independent Poisson
events with rates proportional to abundance, and the habitat prior's
reliability is a free parameter. All of these are configuration knobs or data
swaps.

\paragraph{Conclusion.} Nothing sees the abyss but sound. Sunlight, radar and
radio are gone within metres to a few hundred metres of the surface even for
an ideal detector, so no satellite will ever image a hadal animal; the
satellite's contribution is bathymetry and surface productivity, a prior for
where to look that is worth \SurGainNine\% more species per sensor budget
when it is reliable and nothing when it is not. Low-frequency sound crosses
the full water column --- a 180\,dB call at 50\,Hz is detectable at
\RangeBaleenHyb\,km --- while quiet animals must be within metres of an in-situ
camera or an eDNA sampler. Given sensors in the water, an open-set model can
tell that a call belongs to no catalogued species: distance-based scores on a
learned embedding reach AUROC \NovBestAuroc\ on a physically simulated
soundscape, degrade gracefully with range and catalogue size, and let the
flagged clips be grouped into a handful of candidate species when there are
few of them. On real marine-mammal calls the same pipeline trained from
scratch only reaches \RealBestAuroc, which locates the real bottleneck:
discovery needs an embedding trained on far more bioacoustic data than the
deep sea will ever supply on its own. And whatever the model, the statistics
of rare species set the pace --- the probability that the next detection is
new falls from \SurDiscTen\% to \SurDiscThousand\% between 10 and 1000
sensor-hours. An AI can recognise the unknown; finding it is still a matter
of where you drop the hydrophone, and for how long.
